%%%%%%%%%%  scan idx 616  |  folio 603  %%%%%%%%%%
\uline{section donnée de} \(X\)\uline{), le sous-foncteur de} \(\underline{\Isom}(L,M)\) \uline{formée des isomorphismes compatibles avec ces rigidifications est représentable par un fermé de} \(\underline{\Isom}(L,M)\).
\end{enumerate}

En effet, ii) est clair. Quant à i), rappelons-nous que les foncteurs \(T\mapsto\Hom(L_T,M_T)\) et \(T\mapsto\Hom(M_T,L_T)\) sont représentables par des fibrés vectoriels \(\mathbb V(E)\) et \(\mathbb V(F)\) (d'après EGA III 7.7.8). Alors, si \(u\in\Hom(L_T,M_T)\) et \(v\in\Hom(M_T,L_T)\), les conditions \(u\circ v=1\) et \(v\circ u=1\) définissent un fermé \(F\) de \(\mathbb V(E)\times_s\mathbb V(F)\) et il est clair que \(F\simeq I\).

\textbf{2.5.} Finalement réduisons l'assertion II' à l'assertion suivante:

\(\mathrm{II}^*\). \uline{Le théorème} 1.1 \uline{lorsque} \(X\) \uline{et} \(Y\) \uline{sont intègres sur} \(S\) \uline{et que la suite}
\begin{equation}\tag{*}
  {O}_Y
  \longrightarrow f_*{O}_Y
  \rightrightarrows g_*\bigl({O}_{X\times_yX}\bigr)
\end{equation}
\uline{(où} \(g:X\times_yX\to Y\) \uline{est le morphisme structural) est exacte.}

En fait, il suffit pour ceci de prouver qu'il existe une factorisation de \(f\)
\[
  X\xrightarrow{f_0}Y_0\xrightarrow{f_1}Y_1\longrightarrow\cdots\xrightarrow{f_\eta}Y_\eta=Y
\]
telle que tout \(f_i\) satisfasse à (*) et, quitte à restreindre \(S\), tout \(Y_i\) soit intègre sur \(S\).

Pour ceci, on pose \(Y_0=\Spec(f_*(\underline{O}_X))\), \(f_0\) provenant de la factorisation de Stein de \(f\) (EGA III 4.3.3), de sorte que \(f_{0*}(\underline{O}_X)\xleftarrow{\sim}\underline{O}_{Y_0}\), et que \(f_0\) satisfait donc bien à (*). Pour factoriser de même \(Y_0\to Y\), on utilise le

\textbf{Lemme 2.6.} \uline{Soient} \(Y\) \uline{un schéma noethérien,} \(f:X\to Y\) \uline{un épimorphisme fini (i.e. tel que} \({O}_Y\to f_*({O}_X)\) \uline{soit injectif), alors} \(f\) \uline{se factorise en une suite finie de morphismes finis} \(f_i\) \(X\simeq Y_0\xrightarrow{f_1}Y_1\longrightarrow\cdots\longrightarrow Y_n=Y\) \uline{qui sont des épimorphismes effectifs (i.e. tel que la suite}
